\begin{textsample}{POS Dim 3 – df – Score 14.00 – df/df\_inf\_13.txt}  \label{ex:f3_pos_128}
8.
INSERÇÃO E ACOLHIMENTO É comum tratar sobre a adaptação da \textbf{criança} na etapa Educação \textbf{Infantil}.
Entretanto, não há unanimidade em relação ao termo utilizado para nomear o período de ingresso da \textbf{criança} na \textbf{instituição} de educação para a primeira \textbf{infância}.
Podem ser usados os termos adaptação, acolhimento e inserção.
Como se sabe, a escolha do termo revela concepções sobre as \textbf{crianças} e o modo de condução do trabalho dos profissionais da educação, bem como os pressupostos teóricos que fundamentam a prática educativa.
Nesse sentido, ao tomar como referência os pressupostos teóricos que fundamentam este Currículo, opta -se por outra terminologia, que não adaptação, pois, para a Psicologia Histórico-Cultural, o processo de adaptação contribui para a ocorrência da acomodação, favorecendo a estagnação do desenvolvimento humano, o que não revela a intencionalidade educativa da SEEDF.
Para tal perspectiva teórica, o que contribui para o desenvolvimento humano é o processo de inadaptação, pois esse provoca o ser humano a desenvolver -se. > Se o ambiente não oferecesse nenhum obstáculo a seu desenvolvimento natural, não haveria base alguma para o surgimento de uma ação criadora.
Mas, ao contrário, a inadaptação é a condição principal para o desenvolvimento humano, pois faz emergir a necessidade de transformação do ambiente ( SILVA, 2012, p. 21 ).
Portanto, este Currículo discorre sobre a inserção da \textbf{criança} na Educação \textbf{Infantil} e sobre como precisa ocorrer o seu acolhimento.
Muitas vezes, a inserção da \textbf{criança} em um novo contexto vincula -se às experiências de separação de sua família por um determinado período do \textbf{dia}.
Daí a importância de se debater sobre o acolhimento ( e as formas de efetivá -lo ) como ponto a ser contemplado no planejamento curricular.
Mas por que discutir esse acolhimento na Educação \textbf{Infantil}?
Na verdade, todos os seres humanos vivenciam novas experiências e novos contextos ao longo de sua existência, e, nesse caso, é preciso debater a necessidade de realizar um acolhimento que contribua para o processo de desenvolvimento da capacidade da \textbf{criança} de fazer parte de um novo contexto.
O processo de inserção em novas experiências inicia já com o \textbf{nascimento} da \textbf{criança}, acompanha -a no decorrer de toda sua vida e ressurge a cada nova situação que vivencia.
Como na Educação \textbf{Infantil} se lida com \textbf{bebês}, \textbf{crianças} bem \textbf{pequenas} e \textbf{crianças} \textbf{pequenas} em processo de transição da casa para o mundo mais amplo, o acolhimento ganha ainda mais sentido.
Ressalta -se que esse período pode ser abordado de diferentes pontos de vista : o olhar da \textbf{criança}, das famílias e/ou responsáveis, e o da \textbf{instituição} de educação para a primeira \textbf{infância}.
Ações de acolhimento precisam prever que linguagens, sentimentos, emoções, aprendizagens estejam oportunizando a consolidação da liberdade, da autonomia e do protagonismo \textbf{infantil}, e não apenas respondendo ao cumprimento de ordens com o objetivo de disciplinar os corpos \textbf{infantis} para o modelo escolar tradicional.
Todos, \textbf{crianças} e \textbf{adultos}, são sensíveis ao acolhimento.
Afinal quem não gosta de ser bem recebido?
A qualidade do acolhimento garante o êxito da inserção da \textbf{criança} no contexto da Educação \textbf{Infantil}.
Para que isso ocorra, é fundamental que se faça compreender que o processo de acolhimento exigirá esforços tanto da \textbf{criança} e de seus pais, que buscam adequar -se a essa nova realidade social, como também do professor e \textbf{instituição} educativa, que precisam preparar -se para recebê -la.
Em suma, o estabelecimento de vínculos positivos depende fundamentalmente da forma como a \textbf{criança} e sua família e/ou responsáveis são \textbf{acolhidos} na \textbf{instituição} que oferta Educação \textbf{Infantil}.
O acolhimento da \textbf{criança} envolve aconchego, bem-estar, amparo, \textbf{cuidado} físico e emocional.
Sendo assim, o ato de educar não se separa do ato de \textbf{cuidar}, o que amplia o papel e a responsabilidade dessas \textbf{instituições} nesse momento.
Por isso, a forma como cada uma efetiva o período de acolhida revela a concepção de educação e de \textbf{criança} que orienta suas práticas.
Para tal, o planejamento das atividades é fundamental, para não reproduzir o espontaneísmo e a falta de reflexão.
Pensar como se dará a chegada das \textbf{crianças} ( novas ou não ) nos primeiros \textbf{dias} do \textbf{calendário} e no decorrer do ano letivo, pensar nos tempos, materiais e ambientes, nos profissionais da educação e suas atribuições, nas famílias e/ou responsáveis e suas inseguranças são aspectos importantes para assegurar a qualidade do acolhimento.
Apresentam -se alguns dos aspectos a serem considerados pela \textbf{instituição} educativa no período de acolhimento : - planejamento coletivo ; - envolvimento de todos os profissionais da educação ; - participação das famílias e/ou responsáveis e da comunidade ; - atendimento à diversidade ; - consideração dos sentimentos das \textbf{crianças} e dos \textbf{adultos}.
Dentro do contexto educativo, manifestações, reações e sentimentos podem ser de caráter transitório ou permanente.
Respeitar os jeitos de ser e estar no mundo e os rituais das \textbf{crianças} ajudam em uma transição suave e confiável.
O acolhimento é um princípio a ser concretizado em várias situações que acontecem com as \textbf{crianças} : nos atrasos, no retorno após viagem ou doença, em um acidente ou incidente durante o ano letivo.
Isso porque o acolhimento, para além das datas, materializa a \textbf{humanização} da educação, valendo, portanto, para os primeiros \textbf{dias} e também ao longo do processo educativo.

% matched lemmas: acolher, adulto, bebê, calendário, criança, cuidado, cuidar, dia, humanização, infantil, infância, instituição, nascimento, pequeno
\end{textsample}
